\documentclass[11pt,twoside,a4paper]{article}
% http://www-h.eng.cam.ac.uk/help/tpl/textprocessing/latex_maths+pix/node6.html symboles de math
% http://fr.wikibooks.org/wiki/Programmation_LaTeX Programmation latex (wikibook)
%=========================== En-Tete =================================
%--- Insertion de paquetages (optionnel) ---
\usepackage[french]{babel}   % pour dire que le texte est en fran{\'e}ais
\usepackage{a4}	             % pour la taille   
\usepackage[T1]{fontenc}     % pour les font postscript
\usepackage{epsfig}          % pour gerer les images
%\usepackage{psfig}
\usepackage{amsmath, amsthm} % tres bon mode mathematique
\usepackage{amsfonts,amssymb}% permet la definition des ensembles
\usepackage{float}           % pour le placement des figure
\usepackage{verbatim}

\usepackage{longtable} % pour les tableaux de plusieurs pages

\usepackage[table]{xcolor} % couleur de fond des cellules de tableaux

\usepackage{lastpage}

\usepackage{multirow}

\usepackage{multicol} % pour {\'e}crire dans certaines zones en colonnes : \begin{multicols}{nb colonnes}...\end{multicols} 

% \usepackage[top=1.5cm, bottom=1.5cm, left=1.5cm, right=1.5cm]{geometry}
% gauche, haut, droite, bas, entete, ente2txt, pied, txt2pied
\usepackage{vmargin}
\setmarginsrb{0.20cm}{0.20cm}{0.20cm}{0.20cm}{15pt}{3pt}{15pt}{3pt}

\usepackage{lscape} % changement orientation page
%\usepackage{frbib} % enlever pour obtenir references en anglais
% --- style de page (pour les en-tete) ---
\pagestyle{empty}

% % % en-tete et pieds de page configurables : fancyhdr.sty

% http://www.trustonme.net/didactels/250.html

% http://ww3.ac-poitiers.fr/math/tex/pratique/entete/entete.htm
% http://www.ctan.org/tex-archive/macros/latex/contrib/fancyhdr/fancyhdr.pdf
% \usepackage{fancyhdr}
% \pagestyle{fancy}
% % \newcommand{\chaptermark}[1]{\markboth{#1}{}}
% % \newcommand{\sectionmark}[1]{\markright{\thesection\ #1}}
% \fancyhf{}
% \fancyhead[LE,RO]{\bfseries\thepage}
% \fancyhead[LO]{\bfseries\rightmark}
% \fancyhead[RE]{\bfseries\leftmark}
% \fancyfoot[LE]{\thepage /\pageref{LastPage} \hfill
	% TITLE
% \hfill \includegraphics[width=0.5cm]{img/logo_glider.png} }
% \fancyfoot[RO]{\includegraphics[width=0.5cm]{img/logo_glider.png} \hfill
	% TITLE
% \hfill \thepage /\pageref{LastPage}}
% \renewcommand{\headrulewidth}{0.5pt}
% \renewcommand{\footrulewidth}{0.5pt}
% \addtolength{\headheight}{0.5pt}
% \fancypagestyle{plain}{
	% \fancyhead{}
	% \renewcommand{\headrulewidth}{0pt}
% }


%============================= Corps =================================
\begin{document}

\setlength\parindent{0pt}

\texttt{http://www.capital.fr/carriere-management/dossiers/emploi-les-nouvelles-tendances-du-recrutement-1005275}~\\

\textbf{\LARGE Emploi : les nouvelles m{\'e}thodes de recrutement} ~\\

\emph{\small Mis {\`a} jour le 16/01/15 {\`a} 18:56} ~\\

\textbf{Plus de contacts directs, moins de CV, une traque des candidats plus sophistiqu{\'e}e... Pour embaucher ou pour se faire recruter, des m{\'e}thodes in{\'e}dites apparaissent. En voici sept qui ont le vent en poupe. } ~\\

\begin{minipage}[ht]{7.00cm}
	\includegraphics[width=6.95cm]{img/chercheurs-imagesource_bloc_article_grande_image.jpg}
\end{minipage} \hfill \begin{minipage}[ht]{0.65\textwidth}
	Demain, le recrutement sera mobile. C'est la prochaine r{\'e}volution, apr{\`e}s celle des r{\'e}seaux sociaux professionnels, qui ont d{\'e}j{\`a} boulevers{\'e} les pratiques des candidats et des recruteurs. Les atouts du smartphone ? Sa souplesse et sa r{\'e}activit{\'e}. Or, aujourd'hui, l'objectif des DRH est justement d'optimiser les processus d'embauche. <<La motivation d'un candidat d{\'e}cro{\^i}t au fil du temps pass{\'e} {\`a} s'informer sur une offre>>, observe Jean-Marc Mickeler, associ{\'e} DRH chez Deloitte. D'o{\`u} l'apparition de sites carri{\`e}res Web/ mobile et de plateformes communautaires o{\`u} l'on obtient des r{\'e}ponses plus vite. Oubliez le vieux CV et familiarisez-vous avec ces pratiques originales. %% ~\\
\end{minipage}~\\~\\

\textbf{\textsc{PROSPECTER}}~\\

\textbf{Tendance 1. Piocher l'info directement aupr{\`e}s des salari{\'e}s}~\\

Avoir une vision r{\'e}elle de l'entreprise, de l'int{\'e}rieur, c'est ce que propose le site communautaire Glassdoor, qui vient de d{\'e}barquer des Etats-Unis. Il diffuse des photos des locaux (bureaux, aires de d{\'e}tente, cantines...), des appr{\'e}ciations de salari{\'e}s sur leurs conditions de travail (r{\'e}mun{\'e}rations, avancement, mutuelle...), les notes qu'ils attribuent {\`a} leur entreprise ({\`a} son dynamisme, par exemple) et {\`a} leur patron. Ainsi BNP Paribas a obtenu 3,2 sur 5 et son PDG 70% d'avis favorables. Cet attrait pour l'{\'e}change direct n'a pas {\'e}chapp{\'e} {\`a} certains recruteurs. Depuis 2010, Deloitte affiche sur son site votre carri{\`e}re ( \texttt{deloitterecrute.fr~\footnote{\texttt{http://www.deloitterecrute.fr/}}} ) les contacts d'une centaine d'<<ambassadeurs digitaux>>, des collaborateurs pr{\^e}ts {\`a} r{\'e}pondre {\`a} toute question sur leur v{\'e}cu ou l'int{\'e}r{\^e}t de leurs missions. ~\\

\textbf{Tendance 2. Se faire rep{\'e}rer en usant de canaux inhabituels}~\\

Les Moocs des universit{\'e}s et les Coocs (Corporate open online courses) des entreprises permettent de gagner en visibilit{\'e}. A condition d'y {\^e}tre actif. Les soci{\'e}t{\'e}s partenaires des Moocs rep{\`e}rent sur les plateformes de discussion les participants qui questionnent avec pertinence les enseignants, ceux qui obtiennent les meilleures notes ou qui contribuent aux blogs. Orange, qui lance deux Coocs - d{\'e}coder le code, devenir Web conseiller -, pourrait utiliser ces cours comme canal de recrutement pour ses m{\'e}tiers en tension. ~\\

\textbf{\textsc{POSTULER}}~\\

\textbf{Tendance 3. Envoyer un CV {\`a} partir d'un profil Viadeo ou Linkedin}~\\

Voici une avanc{\'e}e qui va vous faciliter la vie. Vous identifiez, sur Viadeo, une offre d'emploi qui vous int{\'e}resse. H{\'e}las, votre CV n'est pas {\`a} jour. Pas de probl{\`e}me, le r{\'e}seau social r{\'e}cup{\`e}re les informations utiles de votre profil et reconstitue un CV standard que vous pourrez envoyer {\`a} partir de votre mobile. Sur le m{\^e}me principe, certaines entreprises comme colas (travaux publics) et Elior (restauration) proposent sur leur site une fonction Apply with Linkedin (postuler avec Linkedin). ~\\

\textbf{Tendance 4. Vendre votre talent plut{\^o}t que votre exp{\'e}rience}~\\

L'avantage des serious games et autres hackatons, c'est qu'ils confrontent les candidats sur la base de leurs comp{\'e}tences, et non de leur CV. Ainsi, les Battle dev, organis{\'e}es depuis 2013 par RegionsJob, permettent aux d{\'e}veloppeurs informatiques autodidactes, ou au parcours atypique, de d{\'e}fendre leurs chances aupr{\`e}s d'une dizaine de soci{\'e}t{\'e}s sponsors. Chaque participant choisit en amont son langage de programmation (Java, Cobol,.Net...) et les entreprises auxquelles son travail sera soumis. Le concept va {\^e}tre {\'e}tendu bient{\^o}t aux profils commerciaux. ~\\

\textbf{\textsc{EMBAUCHER}}~\\

\textbf{Tendance 5. Diffuser des offres d'emploi gr{\^a}ce {\`a} Twitter}~\\

Se servir de Twitter comme d'une plateforme de recrutement, certaines soci{\'e}t{\'e}s le font d{\'e}j{\`a}. Mais la contrainte des 140 signes laisse peu de place pour mettre en valeur un poste. Un probl{\`e}me r{\'e}solu avec les Twitter Cards con\c{c}ues par Monster. Gr{\^a}ce {\`a} elles, chaque offre d'emploi se transforme automatiquement en un tweet, publi{\'e} sur le fil d'actualit{\'e} de l'entreprise et au contenu enrichi : on peut y ajouter pr{\`e}s de 400 caract{\`e}res (540 en tout) et y ins{\'e}rer des {\'e}l{\'e}ments multim{\'e}dias. ~\\

\textbf{Tendance 6. Traquer les profils convoit{\'e}s}~\\

Pour d{\'e}busquer certains candidats tr{\`e}s demand{\'e}s, comme les d{\'e}veloppeurs et les codeurs - qui ne sont pas pr{\'e}sents sur les r{\'e}seaux sociaux par crainte d'{\^e}tre trop sollicit{\'e}s -, il faut ruser. C'est ce que propose Monster avec TalentBin, un moteur de recherche capable d'identifier les traces num{\'e}riques laiss{\'e}es sur les forums, les blogs et les r{\'e}seaux sociaux d'entraide, et de les agr{\'e}ger. Vous cherchez un d{\'e}veloppeur ma{\^i}trisant un langage particulier ? Soumettez un mot-cl{\'e} (Java, C++, etc.) et les profils les plus pertinents, ainsi que leur contact, appara{\^i}tront. Il ne vous reste plus qu'{\`a} vous montrer convaincant... ~\\

\textbf{Tendance 7. Utiliser un site de cooptation}~\\

C'est une m{\'e}thode efficace pour d{\'e}nicher le mouton {\`a} cinq pattes en un temps record et {\`a} moindres frais. La diff{\'e}rence par rapport {\`a} un cabinet de recrutement classique ? La force du r{\'e}seau. La plateforme de cooptation Keycoopt compte par exemple 15.000 coopteurs en France, r{\'e}partis dans des m{\'e}tiers et des secteurs divers. Chaque annonce co{\^u}te 450 euros {\`a} l'entreprise. Trois candidats, en moyenne, sont s{\'e}lectionn{\'e}s par Keycoopt au terme d'entretiens approfondis. En cas de succ{\`e}s, la soci{\'e}t{\'e} verse {\`a} la plateforme 10% de la r{\'e}mun{\'e}ration annuelle du recrut{\'e}. Quant au coopteur, il touchera 750 euros. ~\\

\emph{Marie-Madeleine S{\`e}ve}~\\

\textbf{Et en 2020 ? Vous serez trait{\'e} comme un client}~\\

A l'avenir, le candidat b{\'e}n{\'e}ficiera d'un feedback en temps r{\'e}el, qui lui {\'e}vitera de fastidieuses relances. Bref, affirme Jean-Christophe Anna, directeur associ{\'e} de Link Humans (recrutement mobile et social), il aura droit aux m{\^e}mes {\'e}gards qu'un client. ~\\

\textbf{JOUR 1} : vous postulez depuis votre mobile. En r{\'e}ponse, un SMS vous invite {\`a} suivre les progr{\`e}s de votre candidature sur une appli de l'entreprise. ~\\

\textbf{J 3} : une alerte sign{\'e}e <<Marc>>, charg{\'e} de recrutement, vous informe que votre profil a plu et vous propose un <<call>>. ~\\

\textbf{J 5} : Marc vous convie {\`a} un test de personnalit{\'e} dans la rubrique <<votre coach>> qui vient enrichir l'appli de l'entreprise. ~\\

\textbf{J 6} : nouvelle alerte : vous avez un entretien dans trois jours et pouvez acc{\'e}der aux r{\'e}sultats de votre test. ~\\

\textbf{J 9} : l'entretien s'est bien pass{\'e}, vous {\^e}tes embauch{\'e}. Vous pouvez d{\'e}buter la pr{\'e}int{\'e}gration via l'appli : visite virtuelle des locaux, infos sur les commodit{\'e}s alentour, livret d'accueil... ~\\
 
\textbf{T{\'e}moignage d'Arnaud Palancade, directeur adjoint du recrutement {\`a} la Soci{\'e}t{\'e} g{\'e}n{\'e}rale}~\\

Pour d{\'e}busquer les perles rares ou des candidats qui ne pensent pas {\`a} nous, nous diffusons nos offres en utilisant le profil Viadeo ou Linkedin de certains collaborateurs. Nous le faisons depuis mai, et cela nous a permis d'{\'e}toffer le volume de candidatures dans la gestion des risques et l'informatique, des d{\'e}partements o{\`u} les postes sont difficiles {\`a} pourvoir. Les experts de ces services deviennent, s'ils le souhaitent, nos ambassadeurs virtuels. Comme ils {\'e}voluent dans l'environnement qui nous int{\'e}resse, les chances d'avoir de bons contacts via leur r{\'e}seau sont {\'e}lev{\'e}es. Il leur suffit d'ajouter sur leur profil trois lignes d{\'e}crivant l'offre. Les annonces sont renouvel{\'e}es tous les quinze jours. Nous avons aujourd'hui 200 ambassadeurs et nous envisageons d'{\'e}largir ce dispositif {\`a} d'autres m{\'e}tiers. ~\\

\textcopyright Management

--- Emploi

\end{document}
